\clearpage
\section{ioctl 要花多少時間}
\label{sec:ioctl}

\begin{oneline}
一次 ioctl 的來回是 \ioctlUs{} µs。一幀 backbone 約 4.17 ms，其中系統呼叫和參數複製只占 \bbSysUs{} µs；看得見的額外成本是 driver 輪詢等硬體：多等 \bbWaitUs{} µs，而且 CPU 每幀花掉 \bbCpuUs{} µs。這個數字是第 3 節改用中斷的理由。
\end{oneline}

\begin{before}[原本怎麼做（第一冊沒有量）]
第一冊說 driver 版本「沒有變慢」，根據的是整條 pipeline 的 fps。但 ioctl 本身要多少時間、driver 等硬體的方式浪費多少 CPU，都沒有量。面試時被問「你的 driver 有多少 overhead」，只能回答「fps 沒掉」，說不出數字。
\end{before}

\begin{tech}
\begin{itemize}
\item \textbf{三個時鐘對照}：app 用 \code{clock_gettime(CLOCK_MONOTONIC)} 量 ioctl 前後；driver 在 ioctl 裡用 \code{ktime_get_ns()} 量 START 到 done，放進回傳的結構；PL 自己的 cycle counter 除以 240 MHz。三者的差就是各層的成本。
\item \textbf{\code{getrusage(RUSAGE_SELF)}}：這個 process 用掉的 CPU 時間（user 加 kernel），量輪詢佔用多少 CPU。
\item \textbf{空 ioctl 迴圈}：\code{TWN_IOC_STATUS} 不拿鎖、只讀六個暫存器，連打 100,000 次取平均。
\end{itemize}
\end{tech}

\begin{gain}
\begin{itemize}
\item ioctl 介面本身不是瓶頸：一次來回 \ioctlUs{} µs，一幀的系統呼叫加參數複製 \bbSysUs{} µs，不到幀時間的 0.1\%。
\item 找到真正的成本：輪詢的等待（driver 量的時間比 PL 的 cycle 多 \bbWaitUs{} µs）與 CPU 時間（每幀 \bbCpuUs{} µs，佔幀時間 \bbCpuPct{}\%）。第 3 節用中斷把它們降到 \irqWaitUs{} µs 與 \irqCpuUs{} µs。
\end{itemize}
\end{gain}

\subsection{實際例子}

\code{test/twn_bench.c} 第 1 部分：空的 ioctl 連打 100,000 次。

\snip{bench_status}

第 2 部分：同一幀反覆跑 1000 次，每次記下三個時間：

\snip{bench_bbrun}

\textbf{板上（開機組 D，driver 輪詢）：}

\boardfile{D2_bench.log}{6-10}

\begin{center}
\small
\begin{tabularx}{\textwidth}{l r X}
\toprule
量的是什麼 & 每幀 & 由哪兩個數字相減 \\
\midrule
app 看到的時間 & 4168.3 µs & ioctl 前後的 \code{clock_gettime} \\
driver 看到的時間 & 4164.8 µs & ioctl 裡 START 前到 done 後的 \code{ktime_get_ns} \\
PL 數的 cycle & 4129.0 µs & 991,000 cycles $\div$ 240 MHz \\
系統呼叫＋複製 & \bbSysUs{} µs & app $-$ driver \\
等待 & \bbWaitUs{} µs & driver $-$ PL \\
這個 process 的 CPU 時間 & \bbCpuUs{} µs & \code{getrusage} 的差除以 1000 \\
\bottomrule
\end{tabularx}
\end{center}

\subsection{原理}

\textbf{一次 ioctl 經過哪些地方。}app 呼叫 \code{ioctl()}，C 函式庫執行 \code{svc} 指令進入 kernel；kernel 的 \code{sys_ioctl} 用檔案描述子找到 \code{struct file}，呼叫 driver 的 \code{unlocked_ioctl}，也就是 \code{twn_ioctl()}；\code{twn_status()} 讀六個暫存器，\code{copy_to_user()} 把 32 bytes 的 \code{struct twn_status} 寫回，再從 kernel 返回。整條路 1.6 µs，其中六次暫存器讀取各要經過 AXI 到 PL 再回來。

\textbf{為什麼「等待」不全是 driver 的錯。}PL 的 cycle counter 只數 17 個 block 執行的時間，block 和 block 之間的空檔不算（第一冊第 8 節）。所以「driver $-$ PL」裡有一部分是 PL 自己的空檔；它是 driver 叫醒延遲的上限，不是精確值。第 3 節換成中斷後，這個差從 \bbWaitUs{} µs 降到 \irqWaitUs{} µs，降下來的部分才確定是輪詢造成的。

\textbf{為什麼輪詢會吃 CPU。}\code{readl_poll_timeout()} 每次睡 13 到 50 µs，用的是高精度計時器（hrtimer）。一幀 4.1 ms 要醒來上百次，每次都要經過計時器中斷、排程器、切回這個 process、讀一次暫存器、再睡。每次只花幾 µs，累積起來每幀約 0.45 ms。

\textbf{量測本身的誤差。}兩次 \code{clock_gettime} 本身也要時間，但它被平均到 100,000 次裡；\code{getrusage} 的解析度是 µs，總量是 0.45 秒等級，誤差可以忽略。同一個開機組裡，「等待」這一項在不同次執行之間會差 10 µs 左右（例如 \code{D2_bench.log} 第 9 行與第 18 行），所以只比較同一次開機、同一個程式的數字。

\subsection{面試會問}

\begin{qa}{你的 driver 有多少 overhead？}
我分三層量。空的 ioctl 一次來回 1.6 µs。跑一幀時，app 量到的時間只比 driver 在 kernel 裡量到的多 3.5 µs，這是系統呼叫加 \code{copy_from_user} 和 \code{copy_to_user}，不到幀時間的 0.1\%。driver 量到的又比 PL 自己數的 cycle 多 30 到 45 µs，這是輪詢的粒度加上 PL 在 block 之間的空檔；而且輪詢讓這個 process 每幀用掉 0.45 ms 的 CPU，大約 11\%。所以介面不是問題，問題是等待的方式。我改成中斷之後，CPU 時間降到每幀 \irqCpuUs{} µs，等待降到 \irqWaitUs{} µs。
\end{qa}

\subsection{關鍵函式速查}

\begin{center}
\small
\begin{tabularx}{\textwidth}{l l X}
\toprule
名稱 & 位置 & 做什麼 \\
\midrule
第 1 部分 & \code{test/twn_bench.c}:\lineof{bench_status} & 空 ioctl 連打 100,000 次 \\
第 2 部分 & \code{test/twn_bench.c}:\lineof{bench_bbrun} & app、driver、PL 三個時間與 CPU 時間 \\
\code{r.ns} & \code{twn_fuse.c} 的 \code{twn_bb_run()} & driver 在 kernel 裡量的 START 到 done \\
\bottomrule
\end{tabularx}
\end{center}
